\subsection{KV Placement Planner}
\label{sec:kv-placement}
\label{sec:capacity-allocation}

Planner (2) protects the hottest $K_l$ sealed blocks in each of the
first $L$ layers, then admits complete footprints within the Peer budget.
\paragraph{Choose protection tiers.}
Let $T_{\mathrm{rc}}$ be historical-prefix recomputation cost,
$T_{\mathrm{rp}}$ known-tail replay cost, and $W_0$ exposed DRAM-loading
wait without Peer residency.
Detection, rerouting, and readmission contribute $R_{\mathrm{floor}}$.
The Planner uses
\begin{equation}
\begin{aligned}
 R_{\mathrm{base}}&=R_{\mathrm{floor}}+T_{\mathrm{rp}},\\
 R_R&=R_{\mathrm{base}}+T_{\mathrm{rc}},\quad
 R_D=R_{\mathrm{base}}+W_0.
\end{aligned}
\label{eq:recovery-base}
\end{equation}
Costs are calibrated by model, context, tail length, hardware, and
representative concurrency, rather than exclusive peak bandwidth.

The three placement tiers are:
\emph{R}, no required backup when $\sigma_r\ge R_R$;
\emph{D}, DRAM only when $R_D\le\sigma_r<R_R$;
and \emph{P}, DRAM coverage plus competition for Peer HBM otherwise.
Tier changes use 20\% hysteresis. Tier-P requests compete by increasing slack; a skipped request retains
its DRAM path. Slack orders admission but does not scale a request's
footprint. Admitting whole footprints makes accounting and updates
predictable under noisy estimates, at the cost of possible unused
capacity or missed gains from partial protection. It is a policy choice,
not an optimality claim or a serial recovery order.

\paragraph{Define protection objects.}
A protection object is a \emph{(layer, block)} containing
all KV heads; the profiled selectors use 64-token blocks.
Sealing makes a full block eligible for copying, whereas the commit
frontier determines what survives failure. Replay forwards the known
uncommitted token IDs to rebuild their KV, including sealed but
uncommitted blocks; it does not resample tokens.

\paragraph{Size offline, rank online.}
The historical blocks needed by a length-$n$ replay form the union
$U_l(n)=\bigcup_{q=1}^{n}\bigcup_g S_{l,q,g}$, where $S_{l,q,g}$ is the
native selector output for query $q$ and group $g$. Repeated reads count
once. A table indexed by model, layer, context bin, and tail length
estimates $h_l=|U_l|/B_l$, giving $K_l=\lceil h_l B_l\rceil$ for $B_l$
historical blocks. Mandatory sink/window blocks must fit \emph{within}
$K_l$; the table is refreshed slowly across requests.

The Adapter counts actual selected indices per layer/block over the
previous 512 decoding steps. The Planner reserves mandatory sealed
blocks, then fills the remaining slots by decreasing count, breaking
ties by recency. An optional block becomes replaceable after two
successive windows outside the top set; replacement reuses its slot
and cannot enlarge $K_l$. Thus calibration sizes the set, while causal
rankings identify its members. Changed members trigger incremental
copies and retain DRAM coverage when removed from Peer.

\paragraph{Protect exposed transfer, not every historical byte.}
The Adapter preloads hot blocks in layer order, skipping Peer-resident
objects. Selector state must arrive before exact demand can be computed;
newly revealed misses take priority over speculative preloads. Cold KV
is filled in the background. Let $\widetilde u$ be the fraction of
historical bytes in the union of preload and demand-fetch sets, and
$h$ the preload fraction, both calibrated from replay traces.
With $N$ homogeneous layers, per-layer replay
compute $C$, and full-layer transfer $D=L_rs/\beta$ ($L_r$ historical
tokens, $s$ KV bytes/token/layer, $\beta$ effective bandwidth),
\begin{align}
 W_{\mathrm{floor}}&=N(\widetilde u-h)D,\label{eq:demand-floor}\\
 W_0&=hD+(N-1)[hD-C]^++W_{\mathrm{floor}}.\label{eq:exposed-wait}
\end{align}
Here $[x]^+=\max(0,x)$. The first two terms are startup preload and
backlog that computation cannot hide. The last is demand wait under a
fixed hot set, causal demand release, and ideal preemptible copy stream.
It is not a lower bound on other hot sets or loading policies.

Peer-resident front-layer hot sets remove the preload component. Set
\begin{equation}
 F^\star(r)=\beta\bigl(W_0-W_{\mathrm{floor}}\bigr),
 \label{eq:peer-byte-target}
\end{equation}
and choose the smallest prefix $L$ with
$\sum_{l=1}^{L}K_l b\ge F^\star(r)$, where $b$ is bytes per layer-block.
Admission also counts selector state. When $hD\le C$, one layer's hot
set suffices. The footprint depends on the request and transfer/compute
profile, not slack, and is recomputed at placement-window boundaries.
For heterogeneous layers this is a planning approximation, not a
runtime guarantee of reaching the homogeneous floor.

\begin{figure}[t]
  \centering
  \includegraphics[width=\columnwidth]{figures/design-single/design-layer-pipeline.pdf}
  \caption{Modeled recovery pipelines: MiniCPM4-8B, 32 layers, 32K
    history, 64-token known tail; $C=2.87$\,ms/layer (measured) and
    $D=5.37$\,ms at an assumed 6.25\,GB/s. In (d), 147\,MiB of Peer
    residency removes preload wait. Totals end at tail-KV reconstruction.
    The homogeneous, preemptible-copy model excludes detection,
    queuing, and interference.}
  \label{fig:design-pipeline}
  \Description{Four timelines compare token-block transfer, full-layer
    pipelining, hot preloads with demand fetch, and Peer-resident hot sets.}
\end{figure}

Figure~\ref{fig:design-pipeline}(a,b) shows that layer pipelining reduces
exposed wait from 172 to 83\,ms; hot preloads in (c) reduce it to
44\,ms; seven Peer-resident hot sets in (d) leave 20\,ms of demand wait.
These calibrated timelines explain placement, not measured fault recovery.

% Continuation of the KV Placement Planner subsection in kv_placement.tex.
\paragraph{Admit whole footprints within the budget.}

Each window admits these footprints using the target KV pool's capacity above
its foreground watermark and admission headroom, capped by
$B_{\mathrm{cap}}$ (default 4\,GiB):
\begin{equation}
 B=[\min(B_{\mathrm{cap}},\mathrm{pool}-\mathrm{watermark}
                         -\mathrm{headroom})]^+.
 \label{eq:peer-budget}
\end{equation}
Algorithm~\ref{alg:placement} admits whole footprints in slack order.
It counts shared physical objects once, skips footprints that do not
fit, and continues to smaller candidates. The set $A_r$ contains the
first $L$ layers' selected sealed blocks and required selector state;
$\mathrm{sealed}(r)$ contains \emph{all} sealed historical KV.

\begin{figure}[t]
  \raggedright
  \refstepcounter{designalgorithm}\label{alg:placement}
  \noindent\textbf{Algorithm \thedesignalgorithm: Whole-footprint placement}
  \par\smallskip
  \noindent Input: requests, slack, profiles, rankings, and budget $B$.
  Output: desired DRAM/Peer sets $D_s,P_s$.
  \par\smallskip
  \begin{tabularx}{\columnwidth}{@{}r@{\quad}X@{}}
    \toprule
    1 & $D_s\gets\varnothing$; $P_s\gets\varnothing$ \\
    2 & \textbf{for} $r$ in $(\sigma_r,\mathrm{ID}_r)$ order: \\
    3 & \quad\textbf{if} $\sigma_r\ge R_R(r)$: \textbf{continue} \\
    4 & \quad$D_s\gets D_s\cup\mathrm{sealed}(r)\cup\mathrm{state}(r)$ \\
    5 & \quad\textbf{if} $\sigma_r\ge R_D(r)$: \textbf{continue} \\
    6 & \quad$A_r\gets\mathrm{footprint}(r,F^\star(r))$ \\
    7 & \quad$\Delta\gets\mathrm{bytes}(A_r\setminus P_s)$ \\
    8 & \quad\textbf{if} $\Delta>B$: \textbf{continue} \\
    9 & \quad$P_s\gets P_s\cup A_r$; $B\gets B-\Delta$ \\
   10 & \textbf{return} $D_s,P_s$ \\
    \bottomrule
  \end{tabularx}
  \Description{Requests retain DRAM coverage unless recomputation fits
    their slack, and gain Peer residency only when their unique footprint fits.}
\end{figure}

Sorting costs $O(n\log n)$ for $n$ requests; each layer-prefix scan costs
$O(N)$ using cached rankings. The Adapter reconciles desired sets with
committed replicas, copying only changes. The target budget is fixed
within a window, but foreground pressure can force reclamation.
Copying, committed, and reclaiming bytes remain physically charged
until release completes; lowering the target does not free addresses.

